\section{Objectifs}

\begin{itemize}
\item Utiliser les notions de multiple, diviseur, nombre pair et nombre impair.
\item Présenter une fraction sous forme irréductible.
\item Manipuler les ensembles : appartenance, inclusion, réunion, intersection, complémentaire, cardinal et produit cartésien.
\item Distinguer implication, réciproque, contraposée, équivalence et contre-exemple.
\item Lire et formuler des propositions avec « et », « ou », négation et quantification en langage naturel.
\item Produire des raisonnements courts par disjonction de cas ou par l'absurde.
\end{itemize}

\section{Prérequis}

\begin{itemize}
\item calcul avec des entiers
\item fractions
\item vocabulaire de base des ensembles
\end{itemize}

\section{Activité de découverte}

On observe les numéros de casiers ouverts après plusieurs passages d'élèves. La question est de savoir quels numéros peuvent être atteints par deux règles différentes : être multiple de 3 et être pair. Cette situation permet de relier arithmétique, ensembles et logique.

\section{Vocabulaire}

\begin{itemize}
\item \textbf{multiple}, \textbf{diviseur}, \textbf{pair}, \textbf{impair}
\item \textbf{ensemble}, \textbf{élément}, \textbf{sous-ensemble}, \textbf{ensemble vide}
\item \textbf{appartenance} \texttt{∈}, \textbf{inclusion} \texttt{⊂}, \textbf{intersection} \texttt{∩}, \textbf{réunion} \texttt{∪}
\item \textbf{complémentaire}, \textbf{cardinal}, \textbf{couple}, \textbf{produit cartésien}
\item \textbf{proposition}, \textbf{implication}, \textbf{réciproque}, \textbf{contraposée}, \textbf{équivalence}, \textbf{contre-exemple}
\end{itemize}

\section{Cours}

\subsection{Arithmétique}

Un entier \texttt{a} est multiple d'un entier \texttt{b} s'il existe un entier \texttt{k} tel que \texttt{a = kb}. Les notions de pair et d'impair se traduisent respectivement par les écritures \texttt{2k} et \texttt{2k+1}.

Deux résultats servent de modèles de démonstration :

\begin{itemize}
\item la somme de deux multiples d'un même entier est encore un multiple de cet entier ;
\item le carré d'un nombre impair est impair.
\end{itemize}

Pour présenter une fraction sous forme irréductible, on simplifie le numérateur et le dénominateur par leurs facteurs communs jusqu'à obtenir deux entiers premiers entre eux.

\subsection{Ensembles}

Un ensemble regroupe des objets appelés éléments. On distingue :

\begin{itemize}
\item \texttt{x ∈ E} : \texttt{x} est un élément de \texttt{E} ;
\item \texttt{A ⊂ E} : \texttt{A} est un sous-ensemble de \texttt{E} ;
\item \texttt{A ∩ B} : éléments communs à \texttt{A} et \texttt{B} ;
\item \texttt{A ∪ B} : éléments appartenant à au moins un des deux ensembles ;
\item \texttt{E \ A} ou le complémentaire de \texttt{A} : éléments de \texttt{E} qui ne sont pas dans \texttt{A} ;
\item \texttt{Card(A)} : nombre d'éléments d'un ensemble fini \texttt{A} ;
\item \texttt{A × B} : ensemble des couples \texttt{(a ; b)} avec \texttt{a ∈ A} et \texttt{b ∈ B}.
\end{itemize}

On utilise aussi l'ensemble vide \texttt{∅}, ainsi que les notations des ensembles de nombres et des intervalles.

\subsection{Logique}

Une proposition mathématique est une phrase à laquelle on peut attribuer une valeur de vérité.

\begin{itemize}
\item « P et Q » n'est vraie que si P et Q sont vraies.
\item « P ou Q » est vraie si au moins l'une des deux propositions est vraie.
\item La négation de « x appartient à A » est « x n'appartient pas à A ».
\item Une implication \texttt{P ⇒ Q} n'est pas automatiquement équivalente à sa réciproque \texttt{Q ⇒ P}.
\item La contraposée de \texttt{P ⇒ Q} est « non Q ⇒ non P » et lui est logiquement équivalente.
\item Un contre-exemple suffit à montrer qu'une affirmation universelle est fausse.
\end{itemize}

En seconde, on lit et on écrit des formulations comme « pour tout réel x » ou « il existe un réel x », sans exiger les symboles \texttt{∀} et \texttt{∃}.

Deux raisonnements sont explicitement travaillés :

\begin{enumerate}
\item \textbf{disjonction de cas} : on sépare une situation en plusieurs cas qui couvrent toutes les possibilités ;
\item \textbf{raisonnement par l'absurde} : on suppose le contraire de ce que l'on veut établir et on aboutit à une contradiction.
\end{enumerate}

\section{Exemples corrigés}

\begin{itemize}
\item \texttt{42} est multiple de \texttt{6} car \texttt{42 = 6 × 7}.
\item Si \texttt{A = {2,4,6,8}} et \texttt{B = {3,6,9}}, alors \texttt{A ∩ B = {6}} et \texttt{A ∪ B = {2,3,4,6,8,9}}.
\item Dans \texttt{E = {1,2,3,4,5}}, si \texttt{A = {2,4}}, alors \texttt{E \ A = {1,3,5}} et \texttt{Card(A)=2}.
\item L'affirmation « si n² est pair alors n est pair » peut être établie par contraposée en montrant que si n est impair, alors n² est impair.
\end{itemize}

\coursefigure{/images/mathematiques/latex-migration/2nde-arithmetique-ensembles-logique-figure-1.svg}{Deux ensembles A et B se recoupent, avec la zone commune mise en évidence.}{Intersection et réunion : les opérations ensemblistes sont ensuite réinvesties en probabilités.}

\section{Algorithmique en Python}

\subsection{Tester un multiple}

\begin{verbatim}
def est_multiple(a, b):
    return b != 0 and a % b == 0

print(est_multiple(42, 6))
\end{verbatim}

L'opérateur \texttt{\%} donne le reste de la division euclidienne. Un reste nul signifie que \texttt{a} est multiple de \texttt{b}.

\section{Erreurs fréquentes}

\begin{itemize}
\item Confondre \texttt{x ∈ A} et \texttt{A ⊂ E}.
\item Confondre une implication et sa réciproque.
\item Employer les symboles \texttt{∀} et \texttt{∃} comme s'ils étaient exigibles : le BO 2026 demande surtout la formulation en langage naturel.
\item Oublier qu'un contre-exemple réfute une affirmation générale mais ne prouve pas une autre affirmation.
\end{itemize}

\section{Formules et idées à retenir}

\begin{itemize}
\item \texttt{a} multiple de \texttt{b} signifie \texttt{a = kb} avec \texttt{k} entier.
\item \texttt{A ∩ B} : éléments communs ; \texttt{A ∪ B} : éléments appartenant à au moins un ensemble.
\item \texttt{Card(A)} désigne le nombre d'éléments d'un ensemble fini.
\item Une implication et sa contraposée sont équivalentes ; la réciproque est une autre proposition.
\end{itemize}

\section{Synthèse}

Le chapitre rassemble l'arithmétique prévue en seconde et les outils de logique qui doivent être réinvestis tout au long de l'année. Les raisonnements, notations et opérations sur les ensembles ne sont donc pas confinés à ce seul chapitre : ils réapparaissent dans les fonctions, la géométrie, les statistiques et les probabilités.
